\documentclass{agentic-ai-template}

\title{ארכיטקטורת סוכני AI}
\author{בר אילן}
\date{מטלה 04 --- קורס AI Agents}

\begin{document}

\maketitlepage

\begin{hebrewabstract}
מסמך זה מסכם את ארכיטקטורת סוכני בינה מלאכותית כפי שנלמדה בהרצאה 07
בקורס \eng{AI Agents}: שלוש שכבות הסוכן (\eng{Software}, \eng{Skill},
\eng{Agent}), מודל שפה גדול וחלון הקשר, \eng{RAG}, כלים וסקילים,
תקשורת בין תוכנות (\eng{SDK}, \eng{GUI}, \eng{CLI}, \eng{API}),
פרוטוקול \eng{MCP}, מגבלות סוכני דפדפן, יצירת \eng{PDF} באמצעות
\eng{LuaLaTeX} ותבניות \eng{CLS}, בקרת איכות (\eng{QA Skills}),
ואורכוסטרציה של ריבוי סוכנים. המסמך כולל פרק יישומי קצר על
\eng{CTI}, \eng{OSINT} וחקירת \eng{Ransomware}. המקור העיקרי לתוכן
הוא תמלול ההרצאה; הביבליוגרפיה משלימה במקורות משניים מאומתים.
\textbf{מילות מפתח:} סוכן \eng{AI}, \eng{RAG}, \eng{MCP}, \eng{LaTeX},
רפטביליות, \eng{CTI}.
\end{hebrewabstract}

\begin{englishabstract}
This document summarizes agentic AI architecture as taught in Lecture 07
of the AI Agents course: the three-layer model (Software, Skill, Agent),
large language models and context windows, retrieval-augmented generation,
tools and skills, inter-application communication (SDK, GUI, CLI, API),
the Model Context Protocol, limitations of browser-based GUI agents,
professional PDF production with LuaLaTeX and custom CLS templates,
QA skills, and multi-agent orchestration. A short applied chapter
covers CTI, OSINT, and ransomware investigation workflows. Primary
source: lecture transcript; bibliography supplements with verified references.
\textbf{Keywords:} AI agent, RAG, MCP, LaTeX, reproducibility, CTI.
\end{englishabstract}

\tableofcontents
\listoffigures
\listoftables
\newpage

\section{מבוא לסוכני AI}

בעשור האחרון עברה הבינה המלאכותית ממודלים של סיווג וחיזוי סטטיסטי, אל מערכות
שמסוגלות לנהל שיחה, לתכנן פעולות ולבצע משימות מורכבות בסביבה דיגיטלית.
במסגרת קורס \eng{AI Agents} נלמד כיצד לבנות, להפעיל ולנהל סוכנים --- ישויות
תוכנה שמשלבות מודל שפה גדול (\eng{LLM}) עם זיכרון עבודה, כלים חיצוניים ויכולת
תכנון. מסמך זה מסכם את הארכיטקטורה כפי שנלמדה בהרצאה 07, ומרחיב אותה לכיוון
יישומי בסייבר.

\subsection{מהו סוכן AI?}

סוכן \eng{AI} הוא מערכת שמקבלת מטרה בשפה טבעית, מפרקת אותה לצעדים, בוחרת כלים
מתאימים ומחזירה תוצאה. בניגוד לצ'אטבוט קלאסי, הסוכן \emph{פועל}: הוא יכול
לקרוא קבצים, להפעיל \eng{API}, לחפש במאגר ידע או לקמפל מסמך \LaTeX{} ל-\eng{PDF}.
לפי חומר ההרצאה, הסוכן אינו עומד בפני עצמו --- הוא נשען על שכבות תחתונות של
תוכנה קבועה ושכבת \eng{Skill} שמעטפת את הקוד בשפה טבעית \cite{brown2020gpt3}.

\begin{insightbox}
המעבר ממודל שפה לסוכן פעיל דורש שלושה מרכיבים: \textbf{תכנון} (מה לעשות),
\textbf{ביצוע} (כלים) ו\textbf{זיכרון} (הקשר + אחסון חיצוני). ללא אחד מהם,
המערכת נשארת מחולל טקסט בלבד.
\end{insightbox}

\subsection{מדוע ארכיטקטורה חשובה?}

ארגונים שמטמיעים סוכני \eng{AI} נתקלים במהירות בבעיות חוזרות: עלויות טוקנים
גבוהות, תוצאות לא עקביות, קושי בבקרת איכות ותלות ב-\eng{GUI} במקום ב-\eng{API}.
הבנת הארכיטקטורה --- שלוש השכבות, \eng{RAG}, \eng{MCP}, תבניות \eng{CLS}
ורב-סוכנים --- מאפשרת לבנות מערכות \textbf{רפטביליות} (\eng{reproducible})
שחוזרות על אותה איכות בכל הרצה.

\subsection{מבנה המסמך}

המסמך בנוי מ-15 פרקים: מיסודות הסוכן ושכבותיו, דרך מנגנוני ידע ותקשורת,
ועד יצירת \eng{PDF} מקצועי, בקרת איכות ואורכוסטרציה. פרק 13 מדגים יישום
קצר בתחום \eng{CTI}/\eng{OSINT} וחקירת \eng{Ransomware}. הפרקים האחרונים
מסכמים את המשמעות הניהולית והטכנולוגית.

\subsection{מטרות למידה}

לאחר קריאת המסמך, הקורא אמור להבין:
\begin{itemize}[rightmargin=2em]
  \item את ההבדל בין שכבת תוכנה, \eng{Skill} וסוכן.
  \item כיצד \eng{RAG} מרחיב את חלון ההקשר מבלי לאמן מחדש את המודל.
  \item מדוע \eng{MCP} שונה מ-\eng{Skill} סטטי עם \eng{API}.
  \item מדוע עבודה טקסטואלית/\eng{API} עדיפה על סוכן דפדפן לייצור מסמכים.
  \item כיצד תבנית \eng{CLS} וסקילי \eng{QA} מבטיחים רפטביליות.
\end{itemize}

\begin{notebox}
מונחים באנגלית מסומנים בפקודה \texttt{\textbackslash eng\{\}} --- לדוגמה:
\eng{Context Window}, \eng{Retrieval-Augmented Generation}.
\end{notebox}

\subsection{הקשר אקדמי}

מטלה 04 בקורס \eng{AI Agents} דורשת הפקת מסמך \eng{PDF} מקצועי באמצעות
\eng{LuaLaTeX} וקובץ \eng{CLS} מותאם. המסמך משקף את יכולת הסטודנט להגדיר
הנחיות לסוכן, לבנות תבנית עיצוב נפרדת ולייצר תוכן אקדמי מלא --- לא שלד ריק.
המקור העיקרי לתוכן הוא תמלול הרצאה 07 (\texttt{sources/transcript-lecture07.txt}).

\section{שלוש שכבות של סוכן: Software, Skill, Agent}

אחת התרומות המרכזיות של ההרצאה היא מודל השכבות השלוש, המסביר כיצד סוכן
\eng{AI} בנוי מרכיבים שונים בעלי אחריות ברורה. המודל מאפשר להבין היכן
להציב לוגיקה קבועה, היכן להשתמש בשפה טבעית ואיפה ``המוח'' של הסוכן פועל.

% Three-layer agent architecture
\begin{figure}[H]
\centering
\begin{tikzpicture}[
  layer/.style={draw=primary, fill=primary!8, rounded corners=4pt,
    minimum width=9cm, minimum height=1.2cm, align=center, font=\small}
]
  \node[layer, fill=accentgreen!15, draw=accentgreen] (agent) at (0,2.4)
    {\textbf{שכבת Agent}\\\eng{LLM + Memory + Context + RAG + Tools + Skills}};
  \node[layer, fill=accentblue!12, draw=accentblue] (skill) at (0,0.8)
    {\textbf{שכבת Skill}\\עטיפת שפה טבעית מעל קוד תבניתי};
  \node[layer, fill=lightgray, draw=secondary] (soft) at (0,-0.8)
    {\textbf{שכבת Software}\\פונקציות ותהליכים קבועים};
  \draw[-{Stealth[length=2.5mm]}, thick, secondary] (soft.north) -- (skill.south);
  \draw[-{Stealth[length=2.5mm]}, thick, secondary] (skill.north) -- (agent.south);
\end{tikzpicture}
\caption{שלוש שכבות הסוכן לפי חומר ההרצאה: תוכנה, סקיל וסוכן}
\label{fig:three-layers}
\end{figure}


\subsection{שכבת Software --- התשתית הקבועה}

בשכבה התחתונה נמצאת התוכנה התבניתית: פונקציות, אלגוריתמים ותהליכים שחייבים
לרוץ באופן דטרמיניסטי. דוגמאות: חישוב מרחק קוסינוס בין וקטורים, קריאת קובץ
\eng{CSV}, שליחת בקשת \eng{HTTP} או קומפילציה של \LaTeX{} ל-\eng{PDF}.
שכבה זו אינה ``מבינה'' שפה --- היא מבצעת הוראות מדויקות.

\subsection{שכבת Skill --- שפה טבעית מעל קוד}

שכבת ה-\eng{Skill} עוטפת את התוכנה במעטפת של שפה טבעית. במקום לספק
פרמטרים מדויקים לפונקציה, המשתמש או הסוכן מתארים בכוונה כללית מה נדרש,
וה-\eng{LLM} בשיתוף ה-\eng{Skill} ממפה את התיאור לפרמטרים הנכונים.
לדוגמה: ``חפש דירות זולות ברחוב רוטשילד'' במקום קואורדינטות מדויקות של
שדות בטופס.

\begin{insightbox}
\eng{Skill} הוא לא ``סוכן'' --- הוא מנגנון שמאפשר לדבר עם תוכנה קיימת
בשפה חופשית. הסוכן משתמש ב-\eng{Skills} כאשר הוא מזהה שמשימה מתאימה
לתיאור של אחד מהם.
\end{insightbox}

\subsection{שכבת Agent --- המוח וההחלטות}

שכבת הסוכן כוללת את ה-\eng{LLM}, את חלון ההקשר (\eng{Context Window}),
את הזיכרון, את ה-\eng{RAG}, את הכלים (\eng{Tools}) ואת היכולת לבחור
\eng{Skills} מתאימים. הסוכן מקבל מטרה, מתכנן, מפעיל כלים ומסכם תוצאות.
בניגוד ל-\eng{Skill} בודד, הסוכן יכול לשרשר מספר פעולות ולהתאים את עצמו
למשוב.

\agenttable{
\caption{השוואת שלוש השכבות}
\label{tab:three-layers}
\begin{tabular}{@{}p{2.2cm}p{3.5cm}p{3.5cm}p{3.5cm}@{}}
\toprule
\textbf{מאפיין} & \textbf{Software} & \textbf{Skill} & \textbf{Agent} \\
\midrule
שפה טבעית & לא & כן (ממשק) & כן (ליבה) \\
קביעות & גבוהה & בינונית & נמוכה \\
תכנון רב-שלבי & לא & מוגבל & כן \\
דוגמה & \eng{cosine()} & תרגום מסמך & חקירת אירוע \\
\bottomrule
\end{tabular}
}

\subsection{זרימת עבודה אופיינית}

כאשר משתמש שואל ``מה בירת מדינת ישראל?'', הסוכן:
\begin{enumerate}[rightmargin=2em]
  \item מנתח את הפרומפט ב-\eng{LLM}.
  \item מחפש \eng{Skill} או כלי חיפוש רלוונטי.
  \item מפעיל את שכבת ה-\eng{Software} (למשל חיפוש ב-\eng{RAG}).
  \item מרכיב תשובה מתוך חלון ההקשר ומחזיר למשתמש.
\end{enumerate}

\begin{warningbox}
בלבול בין שכבות --- למשל הצבת לוגיקה עסקית קריטית רק בפרומפט ללא
קוד תבניתי --- עלול לגרום לתוצאות לא עקביות. כלל אצבע: מה שחייב
להיות זהה בכל הרצה שייך לשכבת \eng{Software} או ל-\eng{CLS}.
\end{warningbox}

\subsection{יישום ארגוני}

בארגון, ניתן לייצר ספריית \eng{Skills} לכל מחלקה: כספים, משאבי אנוש,
\eng{CTI}. הסוכן המרכזי בוחר את ה-\eng{Skill} המתאים לפי תיאור בראש
הקובץ. כך נשמרת אחידות בלי לשכפל קוד בכל סוכן.

\section{LLM, Context Window וזיכרון}

מודל השפה הגדול (\eng{LLM}) הוא ליבת הסוכן, אך חשוב להבין את מגבלותיו:
למודל אין זיכרון פנימי מתמשך בין הרצות, וכל המידע שהוא ``רואה'' מוגבל
לגודל חלון ההקשר (\eng{Context Window}) \cite{brown2020gpt3}.

\subsection{מהו LLM?}

\eng{LLM} (\eng{Large Language Model}) הוא רשת נוירונים מאומנת על כמויות
עצומות של טקסט, המסוגלת לחזות המשך סביר לרצף מילים. בפועל, המודל מבצע
משימות כמו סיכום, תרגום, כתיבת קוד ותכנון --- אך תמיד בתוך מסגרת
הטוקנים שהוזנו אליו בבקשה הנוכחית.

\subsection{Context Window --- זיכרון העבודה}

חלון ההקשר הוא ``שולחן העבודה'' של הסוכן בזמן שיחה. בו מתנהלים:
\begin{itemize}[rightmargin=2em]
  \item \eng{System Prompt} --- ההנחיות הכלליות לסוכן.
  \item היסטוריית השיחה --- שאלות ותשובות קודמות.
  \item תוצאות כלים --- פלט מחיפוש, הרצת קוד, שליפת \eng{RAG}.
  \item תוכן \eng{Skills} שנבחרו לבקשה הנוכחית.
\end{itemize}

כאשר החלון מתמלא, יש לקצץ, לסכם או להעביר מידע לאחסון חיצוני. זו אחת
הסיבות לחשיבות \eng{RAG} ולניהול זיכרון ארוך טווח.

\begin{notebox}
לפי ההרצאה: ``אין לו זיכרון בעצמו --- הכל נמצא אצלנו במחשב וכל פעם
אני טוען את ה-\eng{Context Window} מחדש עם כל מה שהיה עד עכשיו.''
\end{notebox}

\subsection{Memory --- סוגי זיכרון}

נהוג להבחין בין:
\begin{description}[rightmargin=2em, style=nextline]
  \item[זיכרון קצר טווח] תוכן חלון ההקשר הנוכחי.
  \item[זיכרון ארוך טווח] מאגר חיצוני --- קבצים, מסד נתונים, \eng{vector store}.
  \item[זיכרון פרוצדורלי] \eng{Skills}, תבניות \eng{CLS} וכללים קבועים.
\end{description}

\subsection{הרכבת פרומפט מלא}

לפני כל קריאה ל-\eng{LLM}, המערכת מרכיבה חבילה:
\[
  \text{Input} = \text{System} + \text{Skills} + \text{RAG Docs} + \text{User Prompt} + \text{Tool Outputs}
\]
המודל מייצר תשובה; התשובה מתווספת לחלון ההקשר; המחזור חוזר.

\subsection{אסטרטגיות לניהול הקשר}

\begin{enumerate}[rightmargin=2em]
  \item \textbf{סיכום תקופתי} --- סוכן מסכם שיחה ארוכה ומחליף היסטוריה בסיכום.
  \item \textbf{שליפה סלקטיבית} --- רק מסמכים רלוונטיים מ-\eng{RAG} נכנסים לחלון.
  \item \textbf{חלוקה למשימות} --- אורכוסטרטור מעביר תת-משימות לסוכנים עם הקשר מצומצם.
\end{enumerate}

\begin{insightbox}
\eng{Chain-of-Thought} \cite{wei2022chain} משפר ביצועים על ידי הרחבת
הטקסט הביניים בחלון ההקשר --- אך מגדיל צריכת טוקנים. איזון בין
איכות לעלות הוא החלטת ארכיטקטורה.
\end{insightbox}

\subsection{השלכות על עלות}

כל טוקן בחלון ההקשר --- קלט ופלט --- עולה כסף במודלים מסחריים.
שליפת מסמכים מיותרים ל-\eng{RAG}, צילומי מסך מדפדפן והיסטוריה ארוכה
מנפחים את העלות. לכן מעבר מ-\eng{GUI} ל-\eng{API}/טקסט הוא גם החלטה
כלכלית.

\section{RAG: Retrieval-Augmented Generation}

\eng{Retrieval-Augmented Generation} (\eng{RAG}) הוא מנגנון שמאפשר לסוכן
לגשת לידע חיצוני מבלי לאמן מחדש את המודל \cite{lewis2020rag}. בהרצאה
הושווה ה-\eng{RAG} ל``דיסק קשיח'': מאגר גדול המסודר וקטורית לחיפוש סמנטי.

% RAG retrieval flow
\begin{figure}[H]
\centering
\begin{tikzpicture}[
  node distance=1.1cm and 1.4cm,
  block/.style={draw=primary, fill=lightgray, rounded corners=3pt,
    minimum width=2.6cm, minimum height=0.9cm, align=center, font=\footnotesize},
  arr/.style={-{Stealth[length=2mm]}, thick, accentblue}
]
  \node[block] (prompt) {שאילתת\\משתמש};
  \node[block, right=of prompt] (embed) {הטמעה\\וקטורית};
  \node[block, right=of embed] (search) {חיפוש\\סמנטי};
  \node[block, below=of search] (docs) {מסמכים\\רלוונטיים};
  \node[block, left=of docs] (ctx) {\eng{Context}\\\eng{Window}};
  \node[block, left=of ctx] (llm) {\eng{LLM}};
  \node[block, below=of llm] (answer) {תשובה};
  \draw[arr] (prompt) -- (embed);
  \draw[arr] (embed) -- (search);
  \draw[arr] (search) -- (docs);
  \draw[arr] (docs) -- (ctx);
  \draw[arr] (prompt.south) |- (ctx.west);
  \draw[arr] (ctx) -- (llm);
  \draw[arr] (llm) -- (answer);
\end{tikzpicture}
\caption{זרימת \eng{RAG}: מהפרומפט לשליפה וקטורית, הזרקה לחלון הקשר ותשובת המודל}
\label{fig:rag-flow}
\end{figure}


\subsection{עקרון הפעולה}

כאשר משתמש שואל שאלה:
\begin{enumerate}[rightmargin=2em]
  \item הפרומפט מומר לוקטור (הטמעה --- \eng{embedding}).
  \item מתבצע חיפוש במאגר הוקטורים באמצעות מרחק קוסינוס (או מטריקה דומה).
  \item המסמכים הקרובים ביותר סמנטית נשלפים מהמאגר.
  \item המסמכים מצורפים לחלון ההקשר יחד עם הפרומפט.
  \item ה-\eng{LLM} מייצר תשובה המבוססת על המידע שנשלף.
\end{enumerate}

דוגמה מההרצאה: שאלה על בירת ישראל $\rightarrow$ חיפוש בוויקיפדיה ובמסמכים
רלוונטיים $\rightarrow$ הזרקה ל-\eng{Context Window} $\rightarrow$ תשובה מדויקת.

\subsection{מדוע לא לאמן מחדש?}

אימון מחדש של \eng{LLM} על מאגר ארגוני הוא יקר, איטי ומסובך לתחזוקה.
\eng{RAG} מאפשר:
\begin{itemize}[rightmargin=2em]
  \item עדכון ידע בזמן אמת (הוספת מסמכים למאגר).
  \item ציטוט מקורות (איזה מסמך שימש בתשובה).
  \item שליטה בהרשאות (מאגרים נפרדים לפי מחלקה).
\end{itemize}

\subsection{דוגמת זרימה (פסאודו-קוד)}

\begin{lstlisting}[language=Python,caption={Pipeline בסיסי של RAG}]
def rag_answer(question, vector_db, llm):
    q_vec = embed(question)
    docs = vector_db.search(q_vec, top_k=5)
    context = "\n".join(docs)
    prompt = f"Context:\n{context}\n\nQuestion: {question}"
    return llm.generate(prompt)
\end{lstlisting}

\begin{warningbox}
\eng{RAG} אינו מבטיח אמיתות --- אם המאגר מכיל מידע שגוי, המודל עלול
להמחיש אותו בביטחון. נדרשת בקרת איכות על המקורות ואימות תשובות.
\end{warningbox}

\subsection{רכיבי מערכת RAG}

\agenttable{
\caption{רכיבי מערכת \eng{RAG} טיפוסית}
\label{tab:rag-components}
\begin{tabular}{@{}p{3cm}p{4.5cm}p{4.5cm}@{}}
\toprule
\textbf{רכיב} & \textbf{תפקיד} & \textbf{דוגמה} \\
\midrule
Ingestion & פיצול ואינדוקס מסמכים & \eng{chunk} של 512 טוקנים \\
Embedding & המרה לוקטור & \eng{text-embedding-3-small} \\
Vector DB & אחסון וחיפוש & \eng{Chroma}, \eng{Pinecone} \\
Retriever & בחירת קטעים & \eng{top-k}, סף דמיון \\
Generator & תשובה סופית & \eng{GPT-4}, \eng{Claude} \\
\bottomrule
\end{tabular}
}

\subsection{קשר לקורס}

במטלות קודמות בקורס (למשל תרגום רב-שלבי) נעשה שימוש בהשוואת וקטורים
למדידת סטייה סמנטית. \eng{RAG} משתמש באותו עקרון --- מרחק וקטורי ---
למטרה אחרת: שליפת ידע במקום הערכת דמיון בין תרגומים.

\subsection{מגבלות}

\begin{itemize}[rightmargin=2em]
  \item איכות התלויה בגודל ואיכות ה-\eng{chunk}.
  \item חיפוש שגוי אם השאלה מנוסחת רחוק מהמסמכים.
  \item עלות הטמעה ואחסון למאגרים גדולים.
  \item צורך באבטחת מידע --- דליפת מסמכים רגישים לפרומפט.
\end{itemize}

\begin{insightbox}
בפרק 13 נראה כיצד \eng{RAG} יכול לתמוך בחקירת \eng{Ransomware} על ידי
אינדוקס דוחות \eng{CTI}, הערות כופר ו-\eng{IOCs} היסטוריים.
\end{insightbox}

\section{Tools, Skills ו-Agent Skills}

סוכן \eng{AI} מבדיל בין \textbf{יכולת ביצוע} (\eng{Tool}) לבין \textbf{מעטפת
ידע והנחיה} (\eng{Skill}). הבנה נכונה של ההבחנה מאפשרת לבנות מערכות מודולריות
שניתן לתחזק ולהרחיב \cite{yao2023react,openai2024agents}.

\subsection{Tools --- כלי ביצוע}

\eng{Tools} הם התוכנות והפונקציות שהסוכן מפעיל: חיפוש באינטרנט, הרצת
\eng{Python}, קריאת קובץ, קומפילציית \LaTeX{}, שליחת \eng{HTTP}. ה-\eng{LLM}
לא ``מריץ'' קוד בעצמו --- הוא מחליט \emph{איזה} כלי להפעיל ועם אילו פרמטרים,
והסביבה מבצעת בפועל.

\subsection{Skills --- קבצי הנחיה}

\eng{Skill} הוא קובץ (לרוב \eng{Markdown}) המגדיר תפקיד, כללים ודוגמאות.
בתחילת כל \eng{Skill} מופיע \eng{description} שמסביר מתי להשתמש בו. הסוכן
סורק את רשימת ה-\eng{Skills}, בוחר את הרלוונטיים לפרומפט וטוען את תוכנם
לחלון ההקשר --- בדומה לבחירת מומחה מספרייה.

\begin{notebox}
במטלה 02 בקורס נבנו \eng{Skills} נפרדים לכל סוכן תרגום
(\texttt{skills/en\_to\_fr.md} וכו'). אותו עיקרון חל על ייצור \eng{PDF},
\eng{QA} וחקירת סייבר.
\end{notebox}

\subsection{Agent Skills --- סקילים ברמת הסוכן}

\eng{Agent Skills} (כמו ב-\eng{Cursor} או \eng{Claude}) הם מפרטים שניתן
לשתף בין פרויקטים: הוראות מערכת, תבניות פרומפט וכללי בטיחות. הם מהווים
\eng{Memory} פרוצדורלי --- ידע שתמיד זמין לסוכן בלי לאחסן אותו בכל שיחה.

\subsection{תהליך בחירת Skill}

לפי ההרצאה, כשמשתמש שואל שאלה:
\begin{enumerate}[rightmargin=2em]
  \item ה-\eng{LLM} סורק תיאורי \eng{Skills} זמינים.
  \item נבחרים \eng{Skills} עם התאמה סמנטית (למשל ``חיפוש באינטרנט'').
  \item תוכן ה-\eng{Skill} נטען ל-\eng{Context Window}.
  \item הסוכן מפעיל \eng{Tools} שה-\eng{Skill} מגדיר או מרמז עליהם.
\end{enumerate}

\subsection{דוגמה: Skill לקומפילציית LaTeX}

\begin{lstlisting}[caption={קטע מתוך Skill לייצור PDF}]
# LaTeX Compile Skill
When the user requests a PDF document:
1. Use agentic-ai-template.cls (do not inline styles in main.tex)
2. Compile with: lualatex -> bibtex -> lualatex x2
3. Fix recurring RTL issues in the .cls file, not per-document
4. Verify TOC and bibliography before delivery
\end{lstlisting}

\subsection{Skills מול Tools --- סיכום}

\agenttable{
\caption{הבחנה בין \eng{Tool} ל-\eng{Skill}}
\label{tab:tools-skills}
\begin{tabular}{@{}p{2.5cm}p{4.5cm}p{4.5cm}@{}}
\toprule
 & \textbf{Tool} & \textbf{Skill} \\
\midrule
מהות & קוד רץ & הנחיות טקסט \\
מיקום & \eng{src/}, \eng{MCP} & \texttt{skills/*.md} \\
שינוי & מפתח תוכנה & מומחה תחום \\
דוגמה & \eng{vector\_compare.py} & \eng{QA-Table.md} \\
\bottomrule
\end{tabular}
}

\begin{insightbox}
\eng{ReAct} \cite{yao2023react} משלב חשיבה (\eng{Reasoning}) ופעולה
(\eng{Acting}) --- המודל כותב מחשבות, בוחר כלי, מקבל תוצאה וחוזר חלילה.
זהו הבסיס לרב-סוכנים מודרניים.
\end{insightbox}

\section{תקשורת בין תוכנות: SDK, GUI, CLI, API}

תוכנות מודרניות נבנות בשכבות. הבנת השכבות חיונית לתכנון סוכנים שצריכים
לדבר עם מערכות קיימות --- בין אם דרך ממשק אנושי או ישירות דרך ממשק מכונה.

% Software communication stack
\begin{figure}[H]
\centering
\begin{tikzpicture}[
  layer/.style={draw=secondary, fill=white, minimum width=8.5cm,
    minimum height=0.85cm, align=center, font=\small}
]
  \node[layer, fill=accentamber!12] (api) at (0,3.2) {\eng{API} --- חשיפה חיצונית לתוכנה};
  \node[layer, fill=accentblue!10] (cli) at (0,2.2) {\eng{CLI} --- ממשק טקסטואלי בטרמינל};
  \node[layer, fill=accentgreen!10] (gui) at (0,1.2) {\eng{GUI} --- ממשק גרפי למשתמש אנושי};
  \node[layer, fill=primary!10] (sdk) at (0,0.2) {\eng{SDK} --- שכבת פיתוח מעל הפונקציות};
  \node[layer, fill=lightgray] (fn) at (0,-0.8) {פונקציות תוכנה --- ליבת הביצוע};
  \foreach \a/\b in {fn/sdk, sdk/gui, sdk/cli, sdk/api}{
    \draw[-{Stealth[length=2mm]}, thick, primary] (\a.north) -- (\b.south);
  }
\end{tikzpicture}
\caption{מחסנית תקשורת בין תוכנות: פונקציות, \eng{SDK}, \eng{GUI}/\eng{CLI} ו-\eng{API}}
\label{fig:software-stack}
\end{figure}


\subsection{שכבת הפונקציות}

בתחתית נמצאות פונקציות התוכנה: חישוב, חיפוש, כתיבה לדיסק. שכבה זו אינה
נגישה ישירות למשתמש הקצה --- היא ליבת הלוגיקה.

\subsection{SDK --- ערכת פיתוח}

\eng{SDK} (\eng{Software Development Kit}) עוטף את הפונקציות בממשק נוח
למפתחים. כל פעולה חיצונית בתוכנה אמורה לעבור דרך ה-\eng{SDK}. כשבונים
\eng{GUI} או \eng{CLI}, הם קוראים ל-\eng{SDK} ולא ישירות לפונקציות גולמיות.

\subsection{GUI --- ממשק גרפי}

\eng{GUI} (\eng{Graphical User Interface}) מיועד לבני אדם: תפריטים, חלונות,
גרירה בעכבר. דוגמאות: PowerPoint, Excel, דפדפן. סוכן שעובד דרך \eng{GUI}
מדמה משתמש אנושי --- לוחץ, מצלם מסך, מקליד --- וזה יקר בטוקנים ובזמן.

\subsection{CLI --- ממשק שורת פקודה}

\eng{CLI} (\eng{Command Line Interface}) הוא תפריט טקסטואלי בטרמינל:
לחיצה על \eng{A} מבצעת פעולה אחת, \eng{B} --- אחרת. פשוט, מהיר וידידותי
לסוכנים. ההרצאה מדגישה מעבר מעבודה ב-\eng{GUI} לעבודה בטרמינל ובטקסט.

\subsection{API --- חשיפה לעולם החיצון}

\eng{API} (\eng{Application Programming Interface}) מאפשר לתוכנה חיצונית
לבצע פעולות מורשות. כשתוכנה א' רוצה לדבר עם תוכנה ב', הן מתקשרות דרך
\eng{API} --- בתנאי שיודעים לתכנת את החיבור.

\begin{lstlisting}[language=Java,caption={דוגמת בקשת API (JSON)}]
POST /api/v1/documents/compile HTTP/1.1
Content-Type: application/json

{
  "source": "main.tex",
  "engine": "lualatex",
  "runs": 3
}
\end{lstlisting}

\agenttable{
\caption{השוואת דרכי גישה לתוכנה}
\label{tab:access-methods}
\begin{tabular}{@{}p{2cm}p{2.5cm}p{3.5cm}p{3.5cm}@{}}
\toprule
\textbf{שיטה} & \textbf{קהל יעד} & \textbf{יתרון} & \textbf{חיסרון לסוכן} \\
\midrule
\eng{GUI} & אדם & אינטואיטיבי & טוקנים, איטיות \\
\eng{CLI} & מפתח/סוכן & מהיר, טקסט & דורש למידה \\
\eng{API} & מערכות & אוטומציה מלאה & דורש תיעוד \\
\eng{SDK} & מפתח פנימי & שליטה מלאה & לא חשוף תמיד \\
\bottomrule
\end{tabular}
}

\subsection{Skill כעטיפה ל-API}

כפי שנלמד בהרצאה, ניתן לעטוף \eng{API} ב-\eng{Skill} ולאפשר ``\eng{AI via API}''
--- שפה חופשית שמתורגמת לקריאות מובנות. כך תוכנה א' ותוכנה ב' יכולות
לתקשר דרך סוכנים בלי שהמפתח יזכור כל פרמטר.

\subsection{כיוון עתידי}

החזון מההרצאה: מערכות הפעלה מבוססות \eng{AI} שבהן המשתמש מדבר בשפה
חופשית והמערכת בוחרת כלים (\eng{Excel}, Python, דפדפן) בשקיפות. למימוש
חזון זה נדרש שרוב התוכנות יחשפו יכולות בטקסט/\eng{API}, לא רק ב-\eng{GUI}.

\section{MCP: Model Context Protocol}

\eng{Model Context Protocol} (\eng{MCP}) הוא תקן שמגדיר כיצד סוכני \eng{AI}
מתחברים לשירותים חיצוניים --- מסדי נתונים, מאגרי קבצים, כלי פיתוח ועוד
\cite{anthropic2024mcp}. ההרצאה מציגה את \eng{MCP} כאלטרנטיבה מודרנית
לעבודה ישירה עם \eng{API} ו-\eng{Skills} סטטיים.

% MCP vs Skill vs API
\begin{figure}[H]
\centering
\begin{tikzpicture}[
  box/.style={draw, rounded corners=3pt, minimum width=2.5cm, minimum height=1.6cm,
    align=center, font=\footnotesize},
  arr/.style={-{Stealth[length=2mm]}, thick}
]
  \node[box, fill=accentblue!12, draw=accentblue] (agent) at (0,0) {\eng{AI Agent}};
  \node[box, fill=lightgray, draw=secondary] (skill) at (-4,-2.2) {\eng{Skill}\\סטטי מקומי};
  \node[box, fill=accentamber!12, draw=accentamber] (api) at (0,-2.2) {\eng{API}\\נקודת קצה};
  \node[box, fill=accentgreen!12, draw=accentgreen] (mcp) at (4,-2.2) {\eng{MCP Server}\\מחובר מרוחק};
  \node[box, fill=primary!8, draw=primary] (sdk) at (0,-4.2) {\eng{SDK / Software}};
  \draw[arr, accentblue] (agent) -- (skill);
  \draw[arr, accentamber] (agent) -- (api);
  \draw[arr, accentgreen] (agent) -- (mcp);
  \draw[arr, secondary] (skill) -- (sdk);
  \draw[arr, secondary] (api) -- (sdk);
  \draw[arr, secondary] (mcp) -- (sdk);
\end{tikzpicture}
\caption{השוואה מושגית: סוכן יכול לגשת ליכולות תוכנה דרך \eng{Skill}, \eng{API} או \eng{MCP}}
\label{fig:mcp-comparison}
\end{figure}


\subsection{מהו MCP?}

\eng{MCP} הוא ``שרת'' שרץ בנקודת קצה --- במחשב המקומי או אצל ספק התוכנה ---
ומציע לסוכן רשימת כלים (\eng{tools}), משאבים (\eng{resources}) ותבניות
פרומפט. הסוכן ``מכיר'' את ה-\eng{MCP} הרשום אצלו (בדומה לרשימת \eng{Skills})
ומפעיל אותו כשהמשימה מתאימה.

\subsection{MCP מול API ישיר}

בעבודה מול \eng{API} ישיר, המפתח חייב לדעת כתובת, אימות, סכמת בקשות
ועדכוני גרסה. ב-\eng{MCP}, ספק התוכנה מטפל בשכבה זו; הסוכן מגלה יכולות
דינמית מהשרת.

\begin{insightbox}
ההבדל המרכזי מההרצאה: \eng{Skill} סטטי דורש התקנה מחדש כשה-\eng{API}
משתנה; \eng{MCP} מחובר ``אונליין'' ומתעדכן אוטומטית מול היצרן.
\end{insightbox}

\subsection{מודל שרת-לקוח}

\eng{MCP} מבוסס על מודל שרת ולקוח, בדומה לגלישה באינטרנט: הסוכן ``גולש''
לשרת ומבצע פעולות מורשות. אין צורך לדעת כל פרט טכני על מיקום התוכנה ---
רק את נקודת הקצה של ה-\eng{MCP}.

\subsection{דוגמת תצורה (מושגית)}

\begin{lstlisting}[caption={דוגמת הגדרת MCP בקובץ JSON}]
{
  "mcpServers": {
    "filesystem": {
      "command": "npx",
      "args": ["-y", "@modelcontextprotocol/server-filesystem", "./docs"]
    },
    "browser": {
      "command": "cursor-ide-browser",
      "args": []
    }
  }
}
\end{lstlisting}

\subsection{MCP בתוך ארכיטקטורת הסוכן}

כפי שהוסבר: הסוכן מכיל פנייה ל-\eng{MCP}. כמו חיפוש \eng{Skills}, הוא
בוחר את השרת המתאים. בתוך ה-\eng{MCP} עדיין קיים \eng{API} שמדבר עם
ה-\eng{SDK} --- אך הסוכן לא חשוף לפרטים אלה.

\agenttable{
\caption{השוואה: \eng{Skill}+\eng{API} מול \eng{MCP}}
\label{tab:mcp-vs-skill}
\begin{tabular}{@{}p{3cm}p{4.5cm}p{4.5cm}@{}}
\toprule
\textbf{מאפיין} & \textbf{Skill + API} & \textbf{MCP} \\
\midrule
מיקום & מקומי, סטטי & שרת מרוחק/מקומי \\
עדכון & ידני & אוטומטי \\
גילוי יכולות & מתוך קובץ Skill & מתוך שרת \\
תקשורת בין סוכנים & שפה חופשית & שפה חופשית \\
\bottomrule
\end{tabular}
}

\subsection{מגבלות ואבטחה}

\begin{itemize}[rightmargin=2em]
  \item הרשאות --- אילו כלי \eng{MCP} מורשים לסוכן?
  \item אימות --- מפתחות \eng{API} וסודות לא בפרומפט.
  \item זמינות --- תלות ברשת ובשרת חיצוני.
  \item תאימות --- לא כל מוצר תומך ב-\eng{MCP} עדיין.
\end{itemize}

\begin{warningbox}
חשיפת \eng{MCP} עם גישה לקבצים מערכת או דפדפן דורשת הגדרת הרשאות
זהירה --- סוכן שגוי עלול למחוק קבצים או לשלוח מידע רגיש.
\end{warningbox}

\subsection{קשר לפרויקט הנוכחי}

ב-\eng{Cursor}, שרתי \eng{MCP} כמו \eng{cursor-ide-browser} ו-\eng{playwright}
מאפשרים לסוכן לפעול בסביבת הפיתוח. מטלה 04 עצמה מדגימה עבודה טקסטואלית
(\LaTeX) --- גישה עדיפה על סימולציית \eng{GUI} לעריכת מסמכים.

\section{מגבולות עבודה מול GUI וסוכני דפדפן}

ההרצאה מציגה הדגמה של סוכן דפדפן (\eng{Claude} עם תוסף, \eng{Nano Browser})
שמבצע חיפוש ביד2 --- ומדגישה את המחיר הכבד של גישה זו. פרק זה מסכם
חסרונות, אלטרנטיבות וכיוון המעבר לעבודה טקסטואלית.

% GUI agent vs API agent
\begin{figure}[H]
\centering
\begin{tikzpicture}
\begin{axis}[
  ybar,
  bar width=18pt,
  ymin=0, ymax=10,
  ylabel={מדד יחסי (גבוה = פחות יעיל)},
  symbolic x coords={GUI Agent, API/Text Agent},
  xtick=data,
  xticklabel style={font=\small},
  legend style={at={(0.5,-0.22)},anchor=north, legend columns=3, font=\footnotesize},
  width=11cm, height=6cm,
  nodes near coords,
  nodes near coords style={font=\tiny}
]
\addplot[fill=accentamber!70] coordinates {(GUI Agent,8) (API/Text Agent,2)};
\addlegendentry{עלות טוקנים}
\addplot[fill=accentblue!60] coordinates {(GUI Agent,9) (API/Text Agent,3)};
\addlegendentry{מורכבות}
\addplot[fill=accentgreen!60] coordinates {(GUI Agent,7) (API/Text Agent,2)};
\addlegendentry{זמן ביצוע}
\end{axis}
\end{tikzpicture}
\caption{סוכן דפדפן מול סוכן טקסט/\eng{API}: עלות, מורכבות וזמן (אילוסטרציה לפי ההרצאה)}
\label{fig:gui-vs-api}
\end{figure}


\subsection{כיצד פועל סוכן דפדפן?}

הסוכן:
\begin{enumerate}[rightmargin=2em]
  \item מקבל מטרה בשפה טבעית (``חפש דירות ברחוב רוטשילד'').
  \item בונה תוכנית עבודה ומבקש אישור.
  \item פותח דפדפן, מצלם מסך (\eng{screenshot}), מנתח תמונה.
  \item מקליד בשדות, לוחץ כפתורים, חוזר עד להשלמה.
\end{enumerate}

כל צילום מסך וניתוח תמונה צורכים טוקנים רבים. כל טעות \eng{UI} (כיוון
כתיבה עברית, חלון מודאלי) מחייבת ניסיון חוזר.

\subsection{חסרונות שזוהו בשיעור}

סטודנטים והמרצה ציינו:
\begin{itemize}[rightmargin=2em]
  \item \textbf{זמן} --- משימה פשוטה נמשכת דקות ארוכות.
  \item \textbf{עלות} --- תשלום לפי טוקנים; צילומי מסך יקרים במיוחד.
  \item \textbf{אבטחה} --- התחברות לגימייל, סיסמאות, אישורי משתמש.
  \item \textbf{RTL} --- אתרים שכותבים משמאל לימין שוברים הקלדה בעברית.
  \item \textbf{שבירות} --- שינוי בעיצוב האתר שובר את הסוכן.
\end{itemize}

\begin{notebox}
המרצה הזכיר כלי ``הפוך על הפוך'' להחלפת כיוון מקלדת עברית-אנגלית ---
פתרון אנקדוטלי, לא ארכיטקטוני.
\end{notebox}

\subsection{מדוע GUI מפריע לעידן הסוכנים?}

החזון מההרצאה: מערכת הפעלה שמבינה שפה חופשית לא צריכה לחקות לחיצות
עכבר. \eng{GUI} תוכנן לבני אדם --- תפריטים, אייקונים, משוב ויזואלי.
סוכן שעובד דרך \eng{GUI} מבזבז משאבים על סימולציה אנושית במקום לגשת
ישירות ל-\eng{API} או לטקסט.

\subsection{השוואה מעשית}

\agenttable{
\caption{סוכן \eng{GUI} מול סוכן \eng{API}/טקסט}
\label{tab:gui-api}
\begin{tabular}{@{}p{3.5cm}p{4cm}p{4cm}@{}}
\toprule
\textbf{מדד} & \textbf{סוכן דפדפן} & \textbf{סוכן טקסט/\eng{API}} \\
\midrule
קלט לסוכן & תמונות מסך & JSON, טקסט, קוד \\
עלות טיפוסית & גבוהה & נמוכה \\
מהירות & איטית & מהירה \\
אמינות & שבירה בעיצוב & יציבה יותר \\
דוגמה במטלה 04 & עריכת Word ב-\eng{GUI} & \eng{LuaLaTeX} + \eng{CLS} \\
\bottomrule
\end{tabular}
}

\subsection{מתי בכל זאת להשתמש ב-GUI Agent?}

יש מקרים שבהם אין \eng{API}:
\begin{itemize}[rightmargin=2em]
  \item אתרים ישנים ללא ממשק מכונה.
  \item טפסים ממשלתיים חד-פעמיים.
  \item משימות חד-פעמיות שלא שוות פיתוח \eng{API}.
\end{itemize}

אך עבור תהליכים \textbf{רפטביליים} --- דוחות, הצעות מחיר, מסמכי קורס ---
יש להעדיף \eng{LaTeX}, \eng{Markdown} ו-\eng{API}.

\begin{insightbox}
המעבר מ-\eng{GUI} לטקסט הוא לא רק טכני --- הוא ארגוני: חברות שמחשיפות
\eng{API} ו-\eng{MCP} מאפשרות אוטומציה; אלה שנשענות רק על ממשק גרפי
נשארות מחוץ למערכת הסוכנים.
\end{insightbox}

\subsection{סיכום פרק}

סוכן דפדפן הוא כלי עוצמתי למשימות אד-הוק, אך לא אדריכלות ברירת מחדל.
למסמכים אקדמיים, דוחות \eng{CTI} ותהליכי ארגון --- עבודה בטקסט עם תבנית
\eng{CLS} עדיפה, זולה וניתנת לבקרת איכות.

\section{LaTeX, LuaLaTeX ויצירת PDF מקצועי}

\LaTeX{} הוא שפת סימון לייצור מסמכים מקצועיים \cite{lamport1994latex}. במסגרת
ההרצאה הוצג ככלי מרכזי לעידן הסוכנים: מסמך שנכתב בטקסט בלבד, ללא גרירת
עכבר, מתאים להפעלה אוטומטית על ידי \eng{AI}.

\subsection{מדוע LaTeX?}

\begin{itemize}[rightmargin=2em]
  \item \textbf{רפטביליות} --- אותו קוד מייצר אותו \eng{PDF} בכל הרצה.
  \item \textbf{הפרדה} --- תוכן (\texttt{.tex}) נפרד מעיצוב (\texttt{.cls}).
  \item \textbf{מסמכים גדולים} --- תוכן עניינים, ביבליוגרפיה, הפניות אוטומטיות.
  \item \textbf{מתמטיקה} --- נוסחאות ברמה ש-\eng{Word} מתקשה להשיג.
  \item \textbf{סוכן-ידידותי} --- פקודות טקסט במקום קואורדינטות עכבר.
\end{itemize}

המרצה הצהיר במפורש: ``טעות פטאלית לעבוד בוורד'' למסמכים רציניים וגדולים.

\subsection{סוגי קבצי מסמך}

ההרצאה סקרה משפחת מסמכי טקסט:
\begin{description}[rightmargin=2em, style=nextline]
  \item[\texttt{.txt}] טקסט שטוח --- ללא עיצוב, מתאים לקוד וסקילים.
  \item[\texttt{.md}] \eng{Markdown} --- כותרות, הדגשות, תמונות; פופולרי לתיעוד.
  \item[\texttt{.tex}] \LaTeX{} --- מסמכים אקדמיים ומקצועיים ל-\eng{PDF}.
  \item[\texttt{.html}] דפי אינטרנט --- ניתן להמיר ל-\LaTeX{} על ידי סוכן.
\end{description}

\subsection{LuaLaTeX --- הקומפיילר}

\eng{LuaLaTeX} הוא מנוע קומפילציה מבוסס \eng{Lua}, התומך ב-\eng{Unicode}
ובפונטים מודרניים (דרך \eng{fontspec}) \cite{polyglossia2023}. לעברית
זה קריטי: \eng{pdfLaTeX} אינו מספיק לשילוב עברית-אנגלית איכותי.

תהליך הקומפילציה:
\begin{enumerate}[rightmargin=2em]
  \item כותבים \texttt{main.tex} + \texttt{.cls} + \texttt{.bib}.
  \item מריצים \texttt{lualatex main.tex}.
  \item מריצים \texttt{bibtex main} (אם יש ביבליוגרפיה).
  \item מריצים \texttt{lualatex} פעמיים נוספות (תוכן עניינים והפניות).
\end{enumerate}

\begin{lstlisting}[caption={דוגמת שימוש בפקודות התבנית}]
\documentclass{agentic-ai-template}
\begin{document}
\maketitlepage
שילוב עברית עם \eng{API} ו-\eng{RAG}.
\begin{notebox}
תיבת מידע מעוצבת מהקובץ \texttt{.cls}
\end{notebox}
\end{document}
\end{lstlisting}

\subsection{מבנה פרויקט לסוכן}

לפי דוגמת המרצה:
\begin{itemize}[rightmargin=2em]
  \item \texttt{main.tex} --- קובץ ראשי, יכול לכלול \texttt{\textbackslash input\{\}}.
  \item \texttt{agentic-ai-template.cls} --- כל העיצוב והפקודות המותאמות.
  \item \texttt{references.bib} --- מקורות בפורמט \eng{BibTeX}.
  \item \texttt{sources/} --- חומר גלם (תמלול, סיכומים).
\end{itemize}

\subsection{יתרונות לעבודה עם סוכן AI}

\begin{insightbox}
סוכן שכותב \LaTeX{} מייצר \eng{diff} קריא, ניתן לגרסה ב-\eng{Git}, וניתן
לתיקון נקודתי. מסמך \eng{Word} שנוצר על ידי סוכן לעיתים ``שובר'' עיצוב
בלתי נראה עד להדפסה.
\end{insightbox}

\subsection{מגבלות}

\begin{itemize}[rightmargin=2em]
  \item עקומת למידה --- שגיאות קומפילציה דורשות הבנה בסיסית.
  \item עברית-\eng{BiDi} --- בעיות RTL/LTR דורשות תבנית מותאמת.
  \item התקנת חבילות --- \eng{MiKTeX}/\eng{TeX Live} נדרשים מקומית.
\end{itemize}

פרקים 10--11 יראו כיצד \eng{CLS} וסקילי \eng{QA} פותרים חלק מהמגבלות
הללו בצורה רפטבילית.

\section{CLS Templates ו-Reproducibility}

קובץ \texttt{.cls} (\eng{document class}) הוא לב העיצוב במסמכי \LaTeX{}.
המרצה הציג אותו כפרויקט עצמאי שמבטיח שכל מסמך שייצור הסוכן ייראה אחיד
--- עקרון הרפטביליות (\eng{Reproducibility}).

\subsection{מהו קובץ CLS?}

\eng{document class} מגדיר:
\begin{itemize}[rightmargin=2em]
  \item גודל עמוד, שוליים, פונטים (עברית ואנגלית).
  \item עיצוב כותרות, כותרת עליונה ותחתונה.
  \item פקודות מותאמות: \texttt{\textbackslash eng}, תיבות מידע, סגנון טבלאות.
  \item כללים קבועים: קופירייט, מספור, צבעים ארגוניים.
\end{itemize}

במקום שהסוכן יחזור על אותם הגדרות עיצוב בכל \texttt{main.tex}, הוא טוען:
\begin{lstlisting}
\documentclass{agentic-ai-template}
\end{lstlisting}

\subsection{CLS כ-Tool}

המרצה תיאר את ה-\eng{CLS} ככלי (\eng{tool}): ניתן להגדיר בו פונקציות
\LaTeX{} חדשות --- למשל נוסחאות במרכז שורה, בדיקת כיוון כתיבה, או
הוספת כותרת עליונה אוטומטית. כל תיקון עיצובי חוזר נכנס ל-\eng{CLS}
ולא נשאר ``תיקון חד-פעמי'' במסמך.

\subsection{לולאת תיקון רפטבילית}

כאשר מתגלה באג ב-\eng{PDF} (טבלה זזה ימינה, עברית הפוכה):
\begin{enumerate}[rightmargin=2em]
  \item מציגים לסוכן את הבעיה --- \textbf{בלי} לתקן מיד.
  \item מבקשים הסבר והצעת תיקון (בקרת איכות).
  \item אם הבעיה בשפה טבעית --- מעדכנים \eng{Skill}.
  \item אם הבעיה טכנית-עיצובית --- מעדכנים \eng{CLS}.
  \item מריצים קומפילציה מחדש; אם עבר --- הבעיה נפתרה לכל המסמכים העתידיים.
\end{enumerate}

\begin{warningbox}
תיקון נקודתי ב-\texttt{main.tex} במקום ב-\texttt{.cls} יגרום לבאג
לחזור במסמך הבא. זהו ניגוד ישיר לרפטביליות.
\end{warningbox}

\agenttable{
\caption{\eng{LaTeX}/סוכן מול \eng{Word} לרפטביליות}
\label{tab:reproducibility}
\begin{tabular}{@{}p{3.5cm}p{4.5cm}p{4.5cm}@{}}
\toprule
\textbf{מאפיין} & \textbf{LaTeX + CLS + Agent} & \textbf{Word + Agent} \\
\midrule
עקביות עיצוב & גבוהה (תבנית) & נמוכה \\
גרסאות & \eng{Git diff} ברור & קובץ בינארי \\
אוטומציה & קומפילציה מלאה & לעיתים ידני \\
עברית מורכבת & דורש \eng{LuaLaTeX} & בעיות נפוצות \\
\bottomrule
\end{tabular}
}

\subsection{יישום ארגוני}

בארגון, מומלץ:
\begin{itemize}[rightmargin=2em]
  \item תבנית \eng{CLS} אחת לכל סוג מסמך (דוח, הצעת מחיר, מאמר).
  \item \eng{PRD} קצר לכל פרויקט --- מגדיר מבנה וקהל יעד.
  \item סקריפט קומפילציה (\texttt{compile.ps1}) --- אותה פקודה לכל אחד.
\end{itemize}

\subsection{קשר למטלה 04}

מטלה זו \emph{מחייבת} קובץ \texttt{agentic-ai-template.cls} אמיתי.
ההערכה כוללת לא רק תוכן אלא גם את היכולת להפריד תוכן מעיצוב --- מיומנות
מרכזית בהגדרת הנחיות לסוכן.

\begin{insightbox}
``הסוכן שלכם יהיה רפטבילי --- הוא לא כל פעם ייצר משהו אחר'' (תמלול
ההרצאה). זהו עקרון ההנדסה שמאחורי מטלה 04.
\end{insightbox}

\section{QA Skills ובקרת איכות}

ייצור \eng{PDF} אוטומטי במסמכים ארוכים (30--50 עמודים) מייצר באגים חוזרים:
טבלאות שבורות, עברית-אנגלית מעורבבת, תמונות חורגות מהשוליים. המרצה
הציג מערכת \eng{QA Skills} היררכית לטיפול בכך --- מודל שניתן ליישם
בכל פרויקט סוכן.

% QA orchestration hierarchy
\begin{figure}[H]
\centering
\begin{tikzpicture}[
  level 1/.style={sibling distance=4.2cm, level distance=1.6cm},
  level 2/.style={sibling distance=2.2cm, level distance=1.4cm},
  every node/.style={draw, rounded corners=2pt, align=center, font=\footnotesize,
    minimum width=2cm, minimum height=0.7cm}
]
  \node[fill=primary!15, draw=primary, font=\bfseries] {QA Super}
    child { node[fill=accentblue!12] {QA Table} }
    child { node[fill=accentgreen!12] {QA Bidi} }
    child { node[fill=accentamber!12] {QA Image} }
    child { node[fill=lightgray] {QA Code} };
\end{tikzpicture}
\caption{היררכיית \eng{QA Skills} לפי דוגמת המרצה: מנהל-על ומומחי תחום}
\label{fig:qa-hierarchy}
\end{figure}


\subsection{פילוסופיית QA}

כאשר מתגלה בעיה:
\begin{enumerate}[rightmargin=2em]
  \item \textbf{אל תתקן סתם} --- בקש הסבר מה קרה.
  \item החלט: האם הפתרון שייך ל-\eng{Skill} (שפה טבעית) או ל-\eng{CLS} (טכני)?
  \item עדכן את השכבה המתאימה.
  \item הרץ מחדש ובדוק שהתיקון כללי.
\end{enumerate}

\subsection{מבנה היררכי}

לפי דוגמת המרצה:
\begin{description}[rightmargin=2em, style=nextline]
  \item[QA Super] ``מנכ''ל'' --- מכיר את כל סוגי הבעיות, מפעיל מומחים.
  \item[QA Table] מחלקת טבלאות --- יישור, עמודות, עברית-אנגלית בטבלה.
  \item[QA Bidi] כיווניות --- RTL/LTR, מספרים, נוסחאות.
  \item[QA Image] תמונות --- גבולות, כיתובים, מסגרות.
  \item[QA Code] קטעי קוד --- רקע, גלישת שורות, התנגשויות.
\end{description}

כל ``מחלקה'' היא \eng{Skill} (או סוכן עם \eng{Skill} ייעודי) עם כללי
זיהוי (\eng{detector}) ותיקון (\eng{fixer}).

\subsection{Detector ו-Fixer}

\begin{itemize}[rightmargin=2em]
  \item \textbf{Detector} --- רץ על \eng{PDF} או על קוד מקור, מדווח בעיות
    (למשל: ``הטבלה זזה ימינה'').
  \item \textbf{Fixer} --- מקבל דיווח, יודע לתקן (למשל: עדכון הגדרות
    עמודה ב-\eng{CLS}).
\end{itemize}

\subsection{דוגמת Skill QA-Table}

\begin{lstlisting}[caption={קטע מ-QA Skill לטבלאות}]
# QA-Table Skill
Detect:
- Tables aligned to wrong margin in RTL documents
- Column order reversed (Hebrew headers)
Fix:
- Apply \begin{RTL} inside table environment
- Use p{} columns with \RaggedLeft
- Update agentic-ai-template.cls @AtBeginEnvironment{table}
Never: patch only the single table in main.tex
\end{lstlisting}

\subsection{יתרונות מערכת QA}

\agenttable{
\caption{רכיבי מערכת \eng{QA} למסמכי \eng{PDF}}
\label{tab:qa-system}
\begin{tabular}{@{}p{3cm}p{4.5cm}p{4.5cm}@{}}
\toprule
\textbf{רכיב} & \textbf{תפקיד} & \textbf{פלט} \\
\midrule
Detector & זיהוי חריגה & דוח בעיות \\
Fixer & תיקון ממוקד & \eng{diff} ב-\eng{CLS}/\eng{Skill} \\
QA Super & תיאום & אישור לפני הגשה \\
Compile & אימות & \eng{PDF} סופי \\
\bottomrule
\end{tabular}
}

\subsection{יישום במטלה 04}

רשימת בדיקות לפני הגשה:
\begin{itemize}[rightmargin=2em]
  \item קומפילציה ללא שגיאות (\texttt{compile.ps1}).
  \item תוכן עניינים, רשימת איורים וטבלאות מעודכנים.
  \item ביבליוגרפיה עם ציטוטים [1], [2] בגוף המסמך.
  \item לפחות 30 עמודים, 5 טבלאות, 5 איורים.
  \item עברית-אנגלית ללא שבירות בולטות.
\end{itemize}

\begin{insightbox}
\eng{QA Skills} הם ``בטיחות בדם נקנית'' --- כל שורה בקובץ ה-\eng{Skill}
מייצגת באג אמיתי מהעבר. עם הזמן הספרייה הופכת יקרת ערך.
\end{insightbox}

\section{Agent Orchestration: סוכן יחיד מול ריבוי סוכנים}

כמשימות מתמערות, נדרשת החלטה ארכיטקטונית: סוכן אחד עם מגוון \eng{Skills}
(``כובעים'') או מספר סוכנים במקביל, כל אחד עם התמחות. ההרצאה דנה
ביתרונות, חסרונות והצורך ב-\textbf{אורכוסטרטור}.

\subsection{סוכן יחיד, ריבוי כובעים}

סוכן אחד בטרמינל יכול לטעון \eng{Skill} שונה לכל משימה --- היום מומחה
לטבלאות, מחר מומחה לתמונות. היתרונות:
\begin{itemize}[rightmargin=2em]
  \item פשטות ניהולית --- מעקב אחרי תהליך אחד.
  \item עומס חישובי נמוך --- רק instance אחד פעיל.
  \item קל לדבג --- יודעים היכן קרתה שגיאה.
\end{itemize}

החיסרון: ביצוע \textbf{סדרתי}. אם בדיקת טבלאות לוקחת דקה ובדיקת תמונות
דקה נוספת --- סה''כ שתי דקות.

\subsection{ריבוי סוכנים, כובע אחד לכל אחד}

פתיחת מספר טרמינלים (או תהליכים) --- כל סוכן עם \eng{Skill} ייעודי.
יתרון מרכזי: \textbf{מקביליות}. טבלאות בעמוד 2 ותמונות בעמוד 1 יכולים
להיבדק בו-זמנית על אותו \texttt{main.tex}, כל עוד הם לא משנים את אותן
שורות.

דוגמה מההרצאה: 50 משימות של דקה כל אחת ---
\begin{itemize}[rightmargin=2em]
  \item סוכן יחיד: 50 דקות.
  \item 50 סוכנים במקביל: כדקה אחת (בתיאוריה).
\end{itemize}

\subsection{חסרונות ריבוי סוכנים}

\begin{itemize}[rightmargin=2em]
  \item מורכבות ניהולית --- קשה לעקוב אחרי 50 תהליכים.
  \item סיכון לולאות וצריכת טוקנים לא מבוקרת.
  \item מגבלת חומרה --- המחשב ``נחנק'' מעומס.
  \item התנגשויות כתיבה --- נעול קבצים ברמת מערכת ההפעלה.
\end{itemize}

\subsection{תפקיד האורכוסטרטור}

כשמפעילים מספר סוכנים, \textbf{חובה} סוכן-על (\eng{Orchestrator}):
\begin{enumerate}[rightmargin=2em]
  \item מחלק משימות לסוכני משנה.
  \item אוסף תוצאות ומרכיב מסמך שלם.
  \item מפקח על תקציב טוקנים וזמן.
  \item מטפל בהתנגשויות (תור כתיבה לקובץ).
\end{enumerate}

\begin{notebox}
במטלה 02 (תרגום רב-שלבי) יושמה אורכוסטרציה: שלושה סוכנים ברצף עם
\texttt{run\_pipeline.py} כמנהל. אותו עיקרון --- עם אפשרות למקביליות
בשלב ה-\eng{QA}.
\end{notebox}

\agenttable{
\caption{השוואת מודלי אורכוסטרציה}
\label{tab:orchestration}
\begin{tabular}{@{}p{3.5cm}p{4cm}p{4cm}@{}}
\toprule
\textbf{מאפיין} & \textbf{סוכן יחיד} & \textbf{ריבוי סוכנים} \\
\midrule
מהירות & איטית & מהירה (מקביל) \\
ניהול & פשוט & מורכב \\
משאבים & נמוכים & גבוהים \\
דיבוג & קל & קשה \\
צורך באורכוסטרטור & אופציונלי & חובה \\
\bottomrule
\end{tabular}
}

\subsection{המלצה מעשית}

המרצה ממליץ על \textbf{תמהיל}: לא 50 סוכנים בבת אחת, ולא סוכן יחיד
לכל דבר. לדוגמה במטלה 04:
\begin{itemize}[rightmargin=2em]
  \item סוכן תוכן --- כותב פרקים לפי תמלול.
  \item סוכן עיצוב --- מתחזק \eng{CLS}.
  \item סוכן \eng{QA} --- מריץ קומפילציה ובודק \eng{PDF}.
\end{itemize}

\begin{insightbox}
שאלה קלאסית למבחן (לפי ההרצאה): זמן ביצוע 50 משימות של דקה ---
סוכן יחיד לעומת 50 סוכנים. הבנת המקביליות היא מפתח להטמעה ארגונית.
\end{insightbox}

\section{יישום בעולם הסייבר: CTI, OSINT וחקירת אירועי Ransomware}

פרק זה מדגים כיצד עקרונות ארכיטקטורת הסוכן --- \eng{RAG}, \eng{Skills},
\eng{Tools}, \eng{MCP} ואורכוסטרציה --- יכולים לתמוך ב-\eng{Cyber Threat
Intelligence} (\eng{CTI}), \eng{OSINT} וחקירת אירועי \eng{Ransomware}.
הדגמה \textbf{מושגית}; אין כאן מערכת \eng{RAG} חיה או סריקת אינטרנט אמיתית.

\subsection{הקשר: Ransomware ו-CTI}

\eng{Ransomware} הוא תוכנת כופר שמצפינה נתונים ודורשת תשלום. חוקר איומים
צריך לאסוף \eng{IOCs} (\eng{Indicators of Compromise}), לנתח הערות כופר,
להצליב עם מאגרי \eng{MITRE ATT\&CK} \cite{mitre2024attack} ולשתף מודיעין
עם הצוות \cite{nist2018cti}.

\subsection{ארכיטקטורת סוכן לחקירה}

\begin{enumerate}[rightmargin=2em]
  \item \textbf{סוכן איסוף} --- \eng{Skill} ל-\eng{OSINT}: קריאת דוחות, לוגים,
    הערות כופר (טקסט). כלים: קריאת קבצים, \eng{API} למאגרי מידע מורשים.
  \item \textbf{סוכן ניתוח} --- \eng{RAG} על מאגר פנימי: דוחות קודמים,
    \eng{IOCs} היסטוריים, טכניקות \eng{ATT\&CK}.
  \item \textbf{סוכן דיווח} --- ייצור \eng{PDF} מקצועי ב-\LaTeX{} מתבנית
    \eng{CLS} ארגונית (כמו במטלה 04).
  \item \textbf{QA Super} --- בודק עקביות טבלאות \eng{IOC}, עברית-אנגלית
    בשמות דומיינים, וקישורים בביבליוגרפיה.
\end{enumerate}

\subsection{תרחיש לדוגמה}

\textbf{אירוע:} זיהוי קובץ חשוד \texttt{invoice\_2026.exe} והודעת כופר
\texttt{README\_DECRYPT.txt}.

הסוכן מבצע:
\begin{itemize}[rightmargin=2em]
  \item חילוץ hash (\eng{SHA-256}), כתובות \eng{IP}, דומיינים מהלוגים.
  \item שליפה מ-\eng{RAG}: ``האם hash דומה הופיע בעבר?''
  \item מיפוי לטכניקות \eng{T1486} (נתונים מוצפנים לכופר) ב-\eng{ATT\&CK}.
  \item הפקת דוח \eng{PDF} עם טבלת \eng{IOCs} וציר זמן.
\end{itemize}

\agenttable{
\caption{מיפוי תפקידי סוכן בחקירת \eng{Ransomware}}
\label{tab:cti-agents}
\begin{tabular}{@{}p{2.8cm}p{3.5cm}p{3.5cm}p{2.5cm}@{}}
\toprule
\textbf{סוכן} & \textbf{Skill} & \textbf{Tool / MCP} & \textbf{פלט} \\
\midrule
איסוף & \eng{OSINT-Collect} & קריאת קבצים, \eng{API} & רשימת \eng{IOC} \\
ניתוח & \eng{CTI-Analyze} & \eng{RAG}, \eng{MITRE} & ממצאים \\
דיווח & \eng{Report-LaTeX} & \eng{LuaLaTeX} & \eng{PDF} \\
QA & \eng{QA-Table} & בדיקת \eng{PDF} & אישור \\
\bottomrule
\end{tabular}
}

\subsection{דוגמת RAG ל-CTI}

\begin{lstlisting}[language=Python,caption={שליפת הקשר לחקירת כופר}]
query = "SHA256 abc123... ransomware family"
docs = cti_vector_db.search(embed(query), top_k=3)
# docs may include: past incident reports, MISP entries
prompt = build_incident_summary(docs, raw_iocs)
report = llm.generate(prompt)
\end{lstlisting}

\subsection{שיקולי אבטחה ואתיקה}

\begin{warningbox}
סוכן \eng{CTI} חייב לפעול במסגרת מדיניות ארגונית: אין סריקת אינטרנט
ללא הרשאה, אין שיתוף מידע רגיש בפרומפטים לענן ציבורי ללא אנונימיזציה,
ורישום מלא (\eng{audit log}) לכל פעולה.
\end{warningbox}

\subsection{מגבלות ההדגמה}

\begin{itemize}[rightmargin=2em]
  \item המאמר לא מפעיל מערכת אמיתית --- הוא ממחיש עיצוב ארכיטקטוני.
  \item מקורות \eng{CTI} דורשים אימות מול מאגרים מורשים (\eng{MISP}, \eng{VirusTotal}).
  \item \eng{OSINT} חייב לעמוד בחוק ובתנאי שימוש של מקורות.
\end{itemize}

\begin{insightbox}
העיקרון המחבר לקורס: כמו שמסמך אקדמי ב-\LaTeX{} הוא רפטבילי, כך דוח
דוח \eng{CTI} מתבנית \eng{CLS} מאפשר לצוותים לשתף פורמט אחיד
בכל אירוע --- עם \eng{RAG} על ידע היסטורי.
\end{insightbox}

\section{סיכום אינטגרטיבי: המשמעות הניהולית והטכנולוגית}

פרק זה מחבר את הנושאים הטכניים למסגרת החלטות ארגונית: מה משתנה כשארגון
מאמץ סוכני \eng{AI} לא רק כצ'אט, אלא כשכבת ביצוע?

\subsection{מעבר מ-GUI ל-API}

המגמה המרכזית מההרצאה: תוכנות יחשפו יכולות בטקסט ו-\eng{API} במקום
להסתמך על \eng{GUI} בלבד. ארגונית:
\begin{itemize}[rightmargin=2em]
  \item צוותי פיתוח --- עדיפות ל-\eng{API} ו-\eng{MCP} על פני אוטומציית מסך.
  \item צוותי תפעול --- הגדרת \eng{Skills} ותבניות מסמך (\eng{CLS}).
  \item הנהלה --- ציפייה לרפטביליות ולא ל``קסם'' חד-פעמי של סוכן.
\end{itemize}

\subsection{שלוש שכבות בארגון}

ניתן למפות את מודל השכבות לתפקידים:
\begin{description}[rightmargin=2em, style=nextline]
  \item[Software] מערכות ליבה, \eng{ERP}, מאגרי נתונים, קומפיילרים.
  \item[Skill] מפרטי עבודה: איך לכתוב דוח, איך לבדוק טבלה, איך לסכם פגישה.
  \item[Agent] נקודת כניסה למשתמש --- מממש מטרות בשפה טבעית.
\end{description}

\subsection{מודל עלויות}

\begin{itemize}[rightmargin=2em]
  \item טוקנים --- חלון הקשר, צילומי מסך, היסטוריה ארוכה.
  \item זמן --- סוכן יחיד לעומת מקביליות עם אורכוסטרטור.
  \item תחזוקה --- \eng{Skills} ו-\eng{CLS} מקטינים עלות תיקון חוזר.
\end{itemize}

\subsection{רפטביליות כיתרון תחרותי}

ארגון שבנה:
\begin{itemize}[rightmargin=2em]
  \item ספריית \eng{Skills} מאושרת,
  \item תבניות \eng{CLS}/\eng{PDF} אחידות,
  \item מערכת \eng{QA} היררכית,
\end{itemize}
יכול לייצר עשרות דוחות, הצעות מחיר או מסמכי רגולציה באותה איכות ---
במרחק פרומפט. מתחרה שעדיין עורך \eng{Word} ידנית נשאר מאחור.

\subsection{סיכונים וניהול סיכונים}

\begin{itemize}[rightmargin=2em]
  \item הזיות מודל --- במיוחד ב-\eng{CTI} ובמשפט; נדרש אימות אנושי.
  \item דליפת מידע --- פרומפטים עם נתונים רגישים.
  \item תלות בספק --- מודל, \eng{MCP}, ענן.
  \item ריבוי סוכנים ללא פיקוח --- לולאות ועלויות.
\end{itemize}

\agenttable{
\caption{מסגרת החלטה ארגונית לסוכן \eng{AI}}
\label{tab:org-framework}
\begin{tabular}{@{}p{3cm}p{4.5cm}p{4.5cm}@{}}
\toprule
\textbf{שאלה} & \textbf{אם כן} & \textbf{אם לא} \\
\midrule
האם חוזר על עצמו? & \eng{CLS} + \eng{Skills} & חד-פעמי ב-\eng{GUI} \\
האם דורש ידע פנימי? & \eng{RAG} מאובטח & פרומפט כללי \\
האם מחובר לתוכנה? & \eng{API}/\eng{MCP} & סוכן דפדפן \\
האם דורש מהירות? & ריבוי סוכנים + אורכוסטרטור & סוכן יחיד \\
\bottomrule
\end{tabular}
}

\subsection{קשר לתפקיד מומחה AI בארגון}

המרצה ציפה שבוגרי הקורס יהיו ``מובילי הטמעת \eng{AI}'' --- לא רק משתמשי
צ'אט, אלא מגדירי תבניות, \eng{Skills}, תהליכי \eng{QA} וארכיטקטורת
אינטגרציה. מטלה 04 מאמנת בדיוק מיומנות זו: הגדרת \eng{PRD}, בניית \eng{CLS},
והפקת תוצר מקצועי.

\begin{insightbox}
המעבר מטכנולוגיה לערך עסקי: סוכן שמייצר \eng{PDF} אחיד ב-30 שניות
מחליף שעות עיצוב --- אם התשתית (\eng{CLS}, \eng{QA}) כבר קיימת.
\end{insightbox}

\section{סיכום ומסקנות}

מסמך זה סקר את ארכיטקטורת סוכני \eng{AI} כפי שנלמדה בקורס \eng{AI Agents},
הרצאה 07, והרחיב אותה ליישום קצר בתחום הסייבר. להלן סיכום המסקנות העיקריות.

\subsection{ממצאים מרכזיים}

\begin{enumerate}[rightmargin=2em]
  \item \textbf{שכבות} --- סוכן בנוי על \eng{Software}, \eng{Skill} ו-\eng{Agent};
    הבחנה זו מנחה היכן למקם לוגיקה קבועה לעומת שפה טבעית.
  \item \textbf{ידע} --- \eng{LLM} אינו זוכר; \eng{Context Window} ו-\eng{RAG}
    מספקים זיכרון עבודה ואחסון חיצוני.
  \item \textbf{חיבוריות} --- \eng{API}, \eng{SDK}, \eng{CLI} ו-\eng{MCP} מאפשרים
    לסוכנים לפעול בתוכנות; \eng{MCP} מציע עדכון דינמי לעומת \eng{Skill} סטטי.
  \item \textbf{ממשק} --- סוכן \eng{GUI} יקר ואיטי; עבודה טקסטואלית ו-\eng{LuaLaTeX}
    עדיפה למסמכים רפטביליים.
  \item \textbf{איכות} --- \eng{CLS} וסקילי \eng{QA} הופכים תיקונים לכלליים
    ולא חד-פעמיים.
  \item \textbf{קנה מידה} --- אורכוסטרציה מאזנת בין סוכן יחיד (פשטות) לריבוי
    סוכנים (מהירות).
\end{enumerate}

\subsection{תרומת מטלה 04}

מטלה 04 מדגימה את כל העקרונות בפרויקט אחד:
\begin{itemize}[rightmargin=2em]
  \item \texttt{PRD.md} --- הגדרת דרישות לסוכן (מטא-מיומנות הקורס).
  \item \texttt{agentic-ai-template.cls} --- רפטביליות עיצובית.
  \item \texttt{main.tex} --- תוכן מלא בעברית עם מונחי \eng{English}.
  \item \texttt{references.bib} --- ביבליוגרפיה בפורמט \eng{IEEE}.
  \item \texttt{compile.ps1} --- תהליך בנייה שחוזר על עצמו.
\end{itemize}

\subsection{כיווני המשך}

\begin{itemize}[rightmargin=2em]
  \item הרחבת ספריית \eng{QA Skills} לפי באגים אמיתיים בקומפילציה.
  \item חיבור \eng{MCP} למאגרי \eng{CTI} מורשים בפרויקט גמר.
  \item ניסוי אורכוסטרציה: סוכן תוכן + סוכן \eng{CLS} + סוכן \eng{QA} במקביל.
  \item מעבר מתמלול שיעור ל-\eng{RAG} פנימי לשאילתות בקורס.
\end{itemize}

\subsection{מילה אחרונה}

עידן סוכני ה-\eng{AI} מחייב לא רק פרומפטים טובים, אלא \textbf{ארכיטקטורה}:
שכבות, ידע, כלים, תבניות ובקרת איכות. מי ששולט בכלי טקסט כמו \LaTeX{}
ובתבנית \eng{CLS} יכול להפוך את הסוכן ממחולל טיוטות למכונת הפקת מסמכים
מקצועיים --- בקורס, בארגון ובחקירות סייבר.

\begin{notebox}
מקור עיקרי: \texttt{sources/transcript-lecture07.txt}. מקורות נוספים
ברשימת הביבליוגרפיה --- יש לאמת כתובות ופרטים לפני הגשה סופית.
\end{notebox}


\bibliographystyle{IEEEtran}
\bibliography{references}

\end{document}
